\section{Introduction}

Election-night totals are not certification. They are preliminary public
claims, often produced before every ballot category has been accepted,
duplicated, adjudicated, or reconciled. A trustworthy public count system must
therefore explain change, not merely detect change. The same one-vote delta can
be a write-in adjudication correction, a late mail batch, a data-entry error, a
missing source file, or an unexplained alteration. These cases require different
evidence and different public conclusions.

V.0 introduced RCOUNT as a reproducible package format for election-count
claims. This paper focuses on the canvass arithmetic layer: how unofficial,
canvassed, recounted, amended, and certified snapshots can be represented as
status-specific summaries, how public events explain transitions, and how batch
manifests and jurisdiction totals constrain the arithmetic.

The key boundary remains the same. RCOUNT does not decide that a canvassing
board acted lawfully. It can show that a declared canvassed total equals the
published precinct or batch summaries; that a correction event exists; that a
late mail batch reconciles accepted, counted, and rejected ballots; and that the
source files still match their hashes. That is a mechanical replay of public
evidence, not legal certification.

